% MT054 / explanation — canonical editable content.
本题有一个必须保留的区别：配套答案把RBF的强表达能力简写成“总能零训练错误”，但无条件的“总能”并不成立。



\begin{itemize}

\item \textbf{A按原解正确。} 已给定线性可分，说明存在把训练点完全分对的线性边界。\(C\to\infty\)是越来越不容忍违例的硬间隔极限，在可分设定下可得到零分类训练错误。

\item \textbf{B错误。} 支持向量正是参与决定边界的点。移动一个点可能改变最佳间隔、其系数以及边界位置，没有“必定不变”的保证。

\item \textbf{C错误。} 欠拟合表示当前模型或表示不够。盲目降维可能进一步丢掉有用信息，不是普遍补救；若能去除噪声有时有帮助，但选项没有给出这个条件。

\item \textbf{D正确。} 标准SVM的目标是凸二次函数、约束为线性，属于凸二次规划。核形式需使用合法半正定核，以保留相应凸性。

\item \textbf{E按原解选，但需写出条件。} 有限个互异输入、正的RBF参数及足够强的数据拟合下，RBF空间能表达把训练标签分开的规则。实际取某个较小\(C\)，模型可能主动接受错误以减小正则代价。

\end{itemize}


E最简单的反例是：两条输入完全相同，一条标\(+1\)、另一条标\(-1\)。任何固定函数对同一个输入只能给同一个预测，不可能同时答对。因此应区分“原卷答案ADE”与“严格的E需限定”，不能靠更换核消除相同输入的冲突标签。


\textbf{出处：} 2023B期中 Q3，原卷第2页（5分）


\textbf{答案依据：} 2023B期中 Q3 随卷／配套参考答案；条件与解释经复核。
